% 
%  chapter4.tex
%  ISELthesis
%  
%  Created by Matilde Pós-de-Mina Pato on 2012/10/09.
%
\chapter{Planeamento e Gestão do Projeto}
\label{cha:planeamento e gestao}
Neste capítulo, descreve-se a estratégia de planeamento e a gestão adotada para o desenvolvimento do presente trabalho. A complexidade que acompanha a execução de um pipeline orquestrado e a implementação de modelos de aprendizagem, como a R-GCN, exige uma organização rigorosa para assegurar o cumprimento dos objetivos estabelecidos.

Aqui, é projetado o cronograma do projeto, estruturado em fases desde a revisão de trabalhos relacionados e estado da arte até à validação final do sistema desenvolvido. São apresentadas as tarefas de maior relevância, o estado atual de desenvolvimento e o planeamento das restantes atividades. Adicionalmente, inclui-se uma secção com análise de riscos e estratégias para os mitigar, garantindo que o projeto resista a futuros desafios.

\section{Análise SWOT}
\label{sec:swot}
Para fundamentar a estratégia de gestão e avaliar a viabilidade deste trabalho, foi realizada uma análise SWOT. Esta abordagem permite identificar os fatores internos e externos que podem influenciar o sucesso do desenvolvimento do sistema proposto, permitindo uma antecipação de riscos e a valorizar os pontos fortes da solução.

\paragraph{Forças (Strengths)}
Os principais pontos fortes do projeto centram-se na qualidade dos dados utilizados na abordagem técnica adotada:
\begin{itemize}
    \item Utilização de fontes oficiais e reconhecidas (DailyMed, Orange Book e Purple Book), garantindo dados fidedignos e clinicamente relevantes;
    \item Adoção de uma arquitetura baseada em Redes Neuronais Convolucionais Relacionais (R-GCN), adequada à representação e aprendizagem em grafos com múltiplos tipos de relações clínicas;
    \item Estruturação do problema através de Grafos de Conhecimento, permitindo representar de forma fiel as entidades biomédicas e as relações entre elas;
    \item Enquadramento do trabalho num domínio atual e relevante, com aplicação direta ao reposicionamento de fármacos.
\end{itemize}

\paragraph{Fraquezas (Weaknesses)}
Foram identificadas limitações que exigem uma monitorização constante durante a execução:
\begin{itemize}
    \item Custo computacional associado ao treino de modelos de grafos com um número elevado de tipos de relação e de nós.
    \item Dependência da qualidade da informação extraída, o que pode influenciar o desempenho e a fiabilidade das associações inferidas.
\end{itemize}

\paragraph{Oportunidades (Opportunities)}
O contexto externo apresenta condições favoráveis que podem potenciar o impacto e a evolução do projeto:
\begin{itemize}
    \item Potencial para apoiar a identificação de novas indicações terapêuticas para fármacos existentes, contribuindo para a redução de custos e tempos de investigação e desenvolvimento.
    \item Possibilidade de expansão futura do sistema através da integração de novas fontes.
    \item Capacidade de servir como ferramenta para investigadores, permitindo um foco maior nas associações com maior probabilidade de sucesso.
\end{itemize}

\paragraph{Ameaças (Threats)}
Existem fatores externos que podem influenciar negativamente o desenvolvimento ou os resultados do projeto:
\begin{itemize}
    \item Complexidade associada à validação clínica das associações recomendadas, que depende da disponibilidade de evidência científica atualizada.
    \item Possíveis alterações nos formatos ou padrões de disponibilização dos dados regulatórios, que podem exigir ajustes no processo de aquisição de informação.
    \item Exigências regulatórias e éticas associadas ao setor farmacêutico, que condicionam a aplicação de recomendações geradas por sistemas de inteligência artificial.
\end{itemize}

\section{Cronograma e Planeamento Temporal}
\label{sec:cronograma}

O planeamento temporal do projeto encontra-se representado através de um diagrama de Gantt, apresentado no final do documento, que sintetiza as principais fases de desenvolvimento, desde o enquadramento inicial até à submissão final.

O trabalho está dividido nas seguintes fases principais:
\begin{itemize}
    \item Definição do problema e proposta inicial, incluindo a clarificação dos objetivos da dissertação e o enquadramento do domínio de estudo;
    \item Revisão da literatura e estado da arte, com análise crítica de abordagens existentes em grafos de conhecimento e aprendizagem automática no domínio biomédico;
    \item Aquisição e preparação dos dados;
    \item Construção do Grafo de Conhecimento, envolvendo a definição de entidades, relações e a organização da informação num grafo multi-relacional;
    \item Implementação da arquitetura R-GCN, destinada à aprendizagem das representações do grafo e à predição de novas associações entre entidades;
    \item Avaliação e análise dos resultados, incluindo a validação das predições geradas e a análise crítica do desempenho do modelo;
\end{itemize}

